Con l’impiego di dispositivi hardware sempre più rapidi e semplici da installare è aumentata considerevolmente la quantità di dati EEG disponibile in rete. Questo ha permesso di lavorare su dataset contenenti dati fortemente eterogenei provenienti da soggetti, sessioni o strumentazioni differenti. Non è un caso, infatti, che algoritmi di deep learning, che richiedono grandi quantità di dati, o di \emph{transfer learning}, che migliorano la performance \emph{cross-subject}, siano in costante miglioramento \cite{hossain2022status}. In questo capitolo vengono descritti il dataset impiegato nel lavoro di tesi, il protocollo con cui è stato raccolto, l’organizzazione delle registrazioni e i limiti di cui è necessario tenere conto nell’interpretazione dei risultati.

\section{Il dataset EEG Motor Movement/Imagery}
Le analisi presentate in questa tesi sono state condotte sull’EEG Motor Movement/Imagery Dataset, raccolto dal gruppo di sviluppo del sistema BCI2000 e distribuito liberamente attraverso l’archivio PhysioNet. Si tratta di una delle raccolte più ampie disponibili per il paradigma di immaginazione motoria: comprende oltre millecinquecento registrazioni provenienti da 109 volontari, acquisite con 64 elettrodi disposti secondo il sistema internazionale 10-10 e campionate a 160 Hz.
Il dataset risponde ad esigenze specifiche e ha permesso di strutturare il progetto sulle due linee caratterizzanti, costituite da un classificatore e da un predittore. La prima caratteristica di rilievo è la numerosità dei soggetti: l'elevato numero di registrazioni eterogenee fornisce informazioni utili per la misurazione delle variazioni tra soggetti e run distinte. La stessa proprietà favorisce risultati statisticamente più significativi nella valutazione del predittore. La seconda caratteristica riguarda la varietà di compiti richiesti ad ogni soggetto: la presenza, per ciascuna misurazione, di run di immaginazione motoria e di esecuzione motoria rafforza la validità degli algoritmi e permette confronti semplici e diretti. Oltre a queste richieste, sono presenti run a riposo e sono proprio queste ad aver permesso lo sviluppo del predittore.
Queste proprietà riflettono i diffusi utilizzi del dataset in ambito scientifico, dove viene regolarmente impiegato come \emph{benchmark} per ottimizzare e validare i propri algoritmi.

\section{Il protocollo sperimentale}
Ogni soggetto viene sottoposto a quattordici registrazioni, chiamate \emph{run}. Le prime due, della durata di un minuto ciascuna, costituiscono la linea di base: al soggetto viene chiesto di rimanere in uno stato di riposo, rispettivamente ad occhi aperti e chiusi. Le dodici run successive, di circa due minuti l’una, sono dedicate ai compiti veri e propri.
I compiti richiesti sono quattro, ciascuno eseguito tre volte nel corso della sessione. Ogni compito consiste nell'esecuzione o nell'immaginazione di un movimento cinestetico in corrispondenza di un determinato segnale luminoso. Nel primo compito la comparsa di un segnale luminoso sul lato sinistro o destro di uno schermo segnala di iniziare la chiusura e l'apertura ripetuta del corrispettivo pugno. Nel secondo compito, con la stessa metodologia, viene richiesto di immaginare tale movimento. Allo stesso modo il terzo e il quarto compito prevedono l'apertura e la chiusura ripetuta di entrambi i pugni o di entrambi i piedi, prima eseguita e poi immaginata, a seconda della posizione del segnale luminoso, in alto o in basso nello schermo.
All’interno di ogni run le fasi di compito e quelle di riposo si alternano. Il bersaglio resta visibile per circa quattro secondi, durante i quali il soggetto esegue o immagina il movimento; scompare poi e segue un intervallo di riposo di durata simile. Ogni singola esecuzione prende il nome di \emph{trial} e costituisce l’unità elementare dell’esperimento.
Come precedentemente accennato, il lavoro presentato in questa tesi si concentra sulle run di immaginazione motoria. Internamente sono state eseguite delle prove per verificare la coerenza dei risultati dei soggetti tra run di immaginazione ed esecuzione motoria. L'esito del confronto non ha alcuna rilevanza fisiologica e non è stato dunque riportato formalmente: è sufficiente confermare che le prestazioni medie di ogni soggetto risultano sensibilmente migliori nelle run di esecuzione motoria piuttosto che in quelle di immaginazione motoria. Le due run di linea di base, in particolare la prima eseguita ad occhi aperti, sono invece impiegate per caratterizzare il ritmo cerebrale del soggetto a riposo da cui derivare il predittore.

\section{Organizzazione dei dati}
Le registrazioni sono distribuite secondo lo standard BIDS (Brain Imaging Data Structure), una convenzione che stabilisce come debbano essere nominate e organizzate le cartelle e i file di un esperimento neuroscientifico. Il vantaggio è che il percorso di un dato non deve essere costruito manualmente: è sufficiente descrivere quale registrazione si desidera, indicando soggetto, compito e numero di run, perché la libreria di lettura individui il file corrispondente. 
Oltre al segnale, ogni file contiene le annotazioni, ossia i marcatori temporali che indicano che cosa stesse facendo il soggetto in ciascun istante. I marcatori sono di tre tipi: uno identifica gli intervalli di riposo, gli altri due le due classi attive della run. Questi ultimi assumono significati differenti a seconda del compito: nelle run dedicate alla mano indicano il pugno sinistro e quello destro, mentre nelle altre indicano entrambi i pugni ed entrambi i piedi.
I dati si articolano quindi su tre livelli. La sessione è l'insieme delle quattordici run e contiene tutti i dati generati da un singolo soggetto. La run è la singola registrazione continua, della durata di circa due minuti (ad eccezione delle prime due). Il trial è una singola ripetizione del compito e dura circa quattro secondi. Infine, è possibile individuare la finestra (metodo delle sliding windows), ovvero il breve segmento di segnale sul quale il sistema formula una decisione, ricavato suddividendo la registrazione in tratti parzialmente sovrapposti. La distinzione fra trial e finestra è essenziale per una valutazione corretta delle prestazioni e verrà spiegata dettagliatamente nel capitolo successivo.
Ogni run contiene trenta annotazioni alternate fra riposo e compito, per un totale di circa quindici ripetizioni suddivise fra i due stati. Distinguendo inoltre un compito dall'altro (sopra/sotto o destra/sinistra), si ricavano circa sette trial per classe di attivazione.

\section{Limiti del dataset}
La quantità di dati disponibile per ciascun soggetto rappresenta la \textbf{limitazione principale}. Poco più di venti ripetizioni per classe costituiscono una base ridotta sulla quale addestrare e validare un classificatore, e impongono l’adozione di metodi in grado di operare in un regime di poche osservazioni e molte dimensioni.
Un secondo limite riguarda la struttura della sessione. Tutte le registrazioni di un soggetto sono state raccolte in un’unica occasione e senza alcun riscontro immediato: durante l’esperimento il soggetto non osserva mai l’effetto della propria attività cerebrale e non ha quindi modo di correggere la propria strategia. Ne derivano due conseguenze. La prima è l’impossibilità di studiare l’evoluzione della prestazione nel tempo e la trasferibilità di un modello da una sessione all’altra, che costituisce uno dei problemi pratici più rilevanti per le interfacce cervello-calcolatore. La seconda, più delicata, è che \textbf{non è possibile distinguere un soggetto realmente incapace di generare un segnale utilizzabile da un soggetto semplicemente non addestrato}: in questo dataset le due condizioni sono indistinguibili per costruzione.
Vi sono infine alcune irregolarità di minore entità. Tre soggetti sono stati registrati a 128 Hz anziché a 160 Hz, circostanza che non impedisce l’elaborazione ma rende leggermente più rumorose le grandezze stimate su ciascuna finestra; un ulteriore soggetto presenta un numero di annotazioni inferiore all’atteso e dispone quindi di un minor numero di ripetizioni. Le anomalie sono state gestite in modo da non dover escludere alcun soggetto dalla sperimentazione.  
